\chapter{Mental Health Screening}
\section*{What Is This Test?} Mental-health screening uses questionnaires and interviews to identify symptoms of anxiety, depression, substance use, trauma, or other concerns.\section*{Alternative Names} Behavioral-health screening; psychosocial screening; psychiatric screening.\section*{Principle} Validated scales measure symptoms and impairment, but diagnosis requires a clinical assessment.\section*{Why This Test Is Ordered} It is used in primary care, emergency care, pregnancy, chronic disease, schools, and before selected treatments.\section*{Primary vs Secondary Test} Screening is preliminary; interview, safety assessment, and diagnostic evaluation are primary follow-up.\section*{How This Test Is Performed} The patient completes confidential questions and discusses symptoms, functioning, medicines, substances, and safety.\section*{What Preparation Is Needed} No preparation. Answer honestly and report immediate danger or self-harm thoughts.\section*{Understanding Results} A positive screen indicates need for evaluation, not a diagnosis; a negative screen does not exclude emerging or concealed symptoms.\section*{Follow-up Tests} Clinical interview, suicide-risk assessment, substance testing, sleep assessment, or targeted medical tests may follow.\section*{References} \begin{enumerate}\item MedlinePlus. Mental Health Screening.\item U.S. Preventive Services Task Force. Mental-health screening recommendations.\end{enumerate}\clearpage\section*{Understanding Results and Follow-up}
The laboratory report should identify the specimen, method, units, reference interval or decision limit, and whether the result is positive, negative, abnormal, or inconclusive. Reference intervals vary with age, sex, pregnancy, specimen type, assay, and laboratory; a result outside a range is not automatically a diagnosis, and a result within a range does not always exclude disease. Discuss unexpected results with the ordering clinician, who can decide whether repeat or confirmatory testing, treatment, or referral is appropriate. Do not change medicines or delay urgent care based on an isolated result.
\section*{Regulatory Status and Authorization Date}This chapter describes a test category, not a specific commercial product. FDA, CDSCO, EU IVDR, and MHRA approval or registration dates therefore cannot be assigned without the assay manufacturer, product name, instrument, and intended use. For laboratory-developed tests, an individual agency approval date may not exist. See the Preface for the interpretation of regulatory status.
\clearpage
